% Generated by Listas/make_statements.py from pratica-lecture-05.md, item 6. Do not edit.
\textbf{Source.} Muñoz-Garcia \& Toro-Gonzalez (2019), \emph{Strategy and Game Theory: Practice Exercises with Answers}, Ch. 6, Exercise 1 ``Infinitely Repeated Prisoner's Dilemma Game'', p.~274 {[}PDF 285{]}. Answer in the source: yes, a worked answer follows the exercise. Reproduced verbatim (statement only; the book's worked answer follows it).



\begin{quote}

Consider the infinitely repeated version of the following prisoner's dilemma game, where C denotes confess and NC represents not confess (Table 6.1):



\textbf{Part (a)} Can players support the cooperative outcome (NC, NC) as a SPNE by using tit-for-tat strategies, whereby players punish deviations from NC in a past period by reverting to the stage-game Nash equilibrium where they both confess (C, C) for just one period, and then returning to cooperation?



\textbf{Part (b)} Consider now that players play tit-for-tat, but reverting to the stage-game Nash equilibrium where they both confess (C, C) for two periods, and then players return to cooperation. Can the cooperative outcome (NC, NC) be supported as a SPNE?



\textbf{Part (c)} Suppose that players use strategies that punish deviation from cooperation by reverting to the stage-game Nash equilibrium for ten periods before returning to cooperation. Compute the threshold discount factor, \(\delta\), above which cooperation can be sustained as a SPNE of the infinitely repeated game.



{\def\LTcaptype{none} % do not increment counter

\begin{longtable}[]{@{}lll@{}}

\toprule\noalign{}

Player 1 ~Player 2 & C & NC \\

\midrule\noalign{}

\endhead

\bottomrule\noalign{}

\endlastfoot

C & 2, 2 & 6, 0 \\

NC & 0, 6 & 3, 3 \\

\end{longtable}

}



\emph{Table 6.1} Prisoner's dilemma game

\end{quote}

